\subsection*{MT-M16. The full arithmetic jet basis has one exact quantitative gate}
The bounded-multiplicity hypothesis in MT-M14 can be removed.
Fix a primitive unitary Hecke function, assume GRH, and use
the canonical $\Xi_\eta,\psi_\eta,Q_\eta$ of MT-M12.
Allow all actual multiplicities $m_\rho$, without a uniform
bound. Write $\rho=1/2+iT$ and set
\[
 P_\rho=
  \prod_{\substack{\nu\in Z,\ \nu\ne\rho\\
              |\operatorname{Im}\nu-T|\le1}}
              |T-\operatorname{Im}\nu|^{m_\nu},
\]
with empty product one. The following conditions are equivalent:
\begin{enumerate}
\item The first nonzero finite-$L$ Taylor coefficients satisfy
$|L^{(m_\rho)}(\rho,\eta)/m_\rho!|
                       \ge c(1+|T|)^{-C}$ uniformly in $\rho$.
\item The local products satisfy
$P_\rho\ge c(1+|T|)^{-C}$ uniformly in $\rho$.
\item The canonical closed convolution quotient has a
Schauder basis with the prescribed normalized Mellin
zero-jet coefficient functionals.
\item It has such an absolute Schauder basis, its closed
ideal is the full zero-jet kernel, and it is topologically
the rapid sequence space on
$\mathcal J=\{(\rho,k):0\le k<m_\rho\}$, weighted by $1+|T|$.
\end{enumerate}
Constants may differ between these equivalent conditions.
Only GRH for this fixed function is used analytically:
$A_D$ supplies it for Dirichlet functions and $M_H$
supplies it conditionally in general. The equivalent
quantitative conditions remain an extra gate; none is
being deduced here from zero location alone.

\begin{proof}
All constants below may depend on the fixed function and
on the indicated fixed strip. NT-Z10 gives both polynomial
strip growth of the finite $L$-function and a local count
including multiplicities. In particular
\[
       m_\rho=O(\log(2+|T|)),\qquad
       \sum_{\rho\in Z}m_\rho(1+|T|)^{-2}<\infty.       \tag{15}
\]
The finitely many low zeros can be included by adjusting
positive constants in every estimate.

\emph{Equivalence of the two quantitative conditions.}
Take the disk centered at $2+iT$ with a zero-free boundary
of radius $R\in[4,5]$. At large height it contains no
poles or trivial zeros. Its total zero multiplicity is
$M=O(\log(2+|T|))$. Remove its zeros by the finite
Blaschke product, with multiplicities, and write $L=GB$.
As proved in MT-M11, $G$ is bounded above by a polynomial
in $|T|$ and below by an inverse polynomial on the
inner half disk; this uses the Euler lower bound at
the center and Harnack's inequality.

The point $\rho$ lies in that inner disk. Its own leading
Blaschke coefficient is
$[R(1-|a|^2)]^{-m_\rho}$ with $a=-3/(2R)$.
Every other factor at $\rho$ is its distance from $\rho$
times a factor bounded above and below by fixed positive
constants. These statements follow directly from
$(z-a_j)/(1-\overline{a_j}z)$ and $|a|\le3/8$.
All other zeros of distance at most one are in the disk.
The distances of the remaining disk zeros are between
one and $2R$. Since there are $M$ factors with multiplicity,
all constant factors together cost at most a fixed power
of $2+|T|$. Therefore
\[
 c(1+|T|)^{-C}P_\rho
   \le\left|\frac{L^{(m_\rho)}(\rho,\eta)}{m_\rho!}\right|
   \le C(1+|T|)^C P_\rho.                              \tag{16}
\]
This proves (1) if and only if (2).

\emph{A radius adapted to the full multiplicity.}
Assume (2), and let $\delta_\rho$ be the lesser of one
and the distance to the nearest distinct zero.
If a distinct zero $\nu$ has this distance less than one,
the product $P_\nu$ contains the factor
$\delta_\rho^{m_\rho}$. Every other factor in $P_\nu$
is at most one. Hence
\[
 \delta_\rho^{m_\rho}\ge P_\nu
                         \ge c(1+|T|)^{-C}.
\]
When $\delta_\rho=1$ the same conclusion is immediate.
Put $r_\rho=\delta_\rho/4$. Equation (15) now gives
\[
                        r_\rho^{-m_\rho}
                                      \le C(1+|T|)^C. \tag{17}
\]
This is the quantitative replacement for the fixed number
of reciprocal-Taylor steps used in MT-M14.

\emph{Minimum modulus without bounded multiplicities.}
The disk proof of (11) works with all multiplicities once
(2) is available. We give the detail that makes the
constants independent of the largest multiplicity.
For a point $s=\sigma+it$ in a fixed strip, choose the
larger disk from MT-M12 and write again $L=GB$.
If that disk has no zeros, the lower bound for $G=L$
already gives the required polynomial lower bound.
If its nearest zero is outside the inner half disk,
every Blaschke factor at $s$ is at least $1/8$.
Otherwise let $\rho_0$ be a nearest disk zero in that
inner half disk. For every other zero,
$|\rho_0-\nu|\le2|s-\nu|$.
The product of the other distances, with multiplicity,
is at least $P_{\rho_0}$: omitted factors of distance
at most one only increase that product, while all
included distances greater than one are at least one.
The normalized Blaschke denominators and radii cost
only a constant to the total multiplicity
$M=O(\log(2+|t|))$. The lower bound for $G$ therefore
gives
\[
 |L(s,\eta)|\ge c_A(1+|t|)^{-C_A}
             \operatorname{dist}(s,Z)^{m_{\rho_0}}
                       \quad\hbox{if }\operatorname{dist}(s,Z)<1,
                                                               \tag{18}
\]
and a polynomial lower bound when that distance is
at least one. The choice of a nearest zero in (18)
is harmless in a tie. Zeros outside the larger disk
have distance greater than one from $s$.

\emph{Uniform Hermite interpolation.}
Put $m=m_\rho$, $z=s-\rho$, and
$G_\rho(s)=\Xi_\eta(s)/z^m$.
On the circle $|z|=r_\rho$, $\rho$ is the nearest zero.
Divide (18) by $|z|^m$ and use the fixed-strip lower
Stirling bound for the gamma factors. This gives
\[
             |G_\rho(s)|\ge
                 c(1+|T|)^{-C}e^{-\kappa|T|}
                            \quad(|z|=r_\rho).
\]
The function $G_\rho$ has no zero in that disk.
Consequently Cauchy's formula bounds the $j$th Taylor
coefficient of $e^{-az^2}/G_\rho(\rho+z)$ by
\[
       C(1+|T|)^C e^{\kappa|T|}r_\rho^{-j}.
                                                               \tag{19}
\]
For every $j<m$, (17) makes this a polynomial bound
times $e^{\kappa|T|}$, with constants independent of $j,m$.
Use exactly the Taylor truncation $P_{\rho,k}$ and
cardinal function (14) of MT-M14, now for all $k<m$.

Their fixed-strip bounds are uniform. When $|z|\ge1$,
the degree of $z^kP_{\rho,k}$ is at most $m-1$, so
\[
 |z^kP_{\rho,k}(z)/z^m|
       \le m\,C(1+|T|)^C e^{\kappa|T|}.
 \]
There is no growing positive power of $|z|$.
The factor $m$ is absorbed using (15), and the Gaussian
bound (8) follows from (5), just as before.
When $|z|<1$, the maximum principle applied to
$G_\rho$ on $|z|=2$ gives
$|G_\rho(s)|\le C(1+|T|)^C e^{-\kappa|T|}2^{-m}$.
The polynomial $z^kP_{\rho,k}$ is bounded by
$m$ times the coefficient bound (19), with (17).
This proves (8) on the remaining region as well,
after increasing a fixed constant.
Inverse transforms $h_{\rho,k}$ therefore satisfy
the polynomial seminorm bounds (9) uniformly over
every jet index $k<m_\rho$.

The normalized jets of any $f\in\mathcal E$ also have
uniform rapid height decay for \emph{all} $k\ge0$:
for each fixed $j$ and each fixed real-part strip,
\[
 \sup_{k\ge0}
 (1+|t|)^j\left|\frac{(\mathcal Mf)^{(k)}(\sigma+it)}{k!}\right|
                    \le C_{A,j}\,p_{N,j}'(f),          \tag{20}
\]
where $p_{N,j}'$ is a finite maximum of defining
seminorms, with sufficiently large fixed $N$.
Indeed, integrate by parts $j$ times in the transform
of $e^{\sigma x}x^k f(x)/k!$.
Every derivative of $x^k/k!$ is either zero or
$x^{k-\ell}/(k-\ell)!$; the remaining coefficients
depend only on $A,j$. For $N-A>1$, its weighted
$L^1$ norm is at most
$2/(N-A)^{k-\ell+1}\le2$, times the corresponding
seminorm of $f$. The region $|t|\le1$ follows from
the same integral bound without integration by parts.
This proves (20).

Equations (15), (20), and the uniform cardinal bounds
give absolute interpolation and its continuous inverse
for every rapid sequence on $\mathcal J$.
Finally, division proves the full ideal equality.
If $F=\mathcal Mf$ has all prescribed zero jets, its
Taylor remainder at a nearest zero is
\[
 F(s)=z^m\int_0^1m(1-u)^{m-1}
                  \frac{F^{(m)}(\rho+uz)}{m!}\,du .
 \]
The measure in this integral has total mass one.
Estimate (20) thus bounds the normalized remainder
uniformly in $m$, and its $|z|^m$ cancels the distance
factor in (18). The quotient $F/L$ is rapidly
decreasing on every fixed strip at large height.
The entire quotient $F/\Xi_\eta$, multiplied by
$e^{\epsilon s^2}$, is the Mellin transform of an
element of $\mathcal E$ for every $\epsilon>0$.
The Gaussian approximate-identity argument of
MT-M12 proves that $f$ lies in the closed ideal.
The reverse inclusion is immediate.
This yields the asserted topological rapid-sequence
identification and (4). Statement (4) implies (3),
and (3) implies (1) by MT-M11. All four conditions
are therefore equivalent.
\end{proof}

This removes both the simplicity assumption of MT-M12 and the
uniform multiplicity assumption of MT-M14. It does not prove the
remaining leading-coefficient or cluster-product bound.
The local $O(\log|T|)$ count is used together with that
quantitative condition; MT-M09 shows why the count alone
cannot replace it.

\subsection*{MT-M17. Nuclear trace with arbitrary zero multiplicities}
Under the equivalent conditions of MT-M16, every
$f\in\mathcal E$ acts by nuclear convolution on $Q_\eta$
and by a trace-class operator on $\ell^2(\mathcal J)$.
Its Hilbert trace is
\[
       \operatorname{Tr}C_f
                  =\sum_{\rho\in Z}m_\rho\mathcal Mf(\rho).
\]
For $f=h*h^\sharp$ this is the nonnegative sum
$\sum_\rho m_\rho|\mathcal Mh(\rho)|^2$.
The faithful-operator positivity obstruction of MT-M15
persists whenever any zero is multiple.

To prove these assertions without a uniform block-size
hypothesis, note that (15) bounds the total multiplicity
in each unit-height interval by $O(\log(2+|T|))$.
The number of Toeplitz matrix units in that interval is
therefore at most $O(\log^2(2+|T|))$.
Estimate (20) makes every matrix entry of $C_f$ uniformly
rapidly decreasing in height, independently of its jet
index. Thus the sum of the trace norms of all matrix
units is finite, giving the trace-class assertion and
the displayed diagonal trace. For nuclearity use the
equicontinuous normalized-jet functionals, the uniform
polynomial seminorm bounds for $h_{\rho,k}$, and the
scalar weight $(1+|T|)^{-2}$ for each matrix unit.
These weights are summable even with the
$O(\log^2(2+|T|))$ count, and the resulting vector
family is bounded in every seminorm by (20).
This is the same explicit nuclear decomposition as
MT-M15, now with all indices included.
The involution formula and the finite local nilpotent
argument in MT-M15 make no bounded-multiplicity
assumption, so they prove the remaining statements.

This is an exact reduction for the prescribed normalized-jet
coordinates. MT-M18 and MT-M19 below supply a different,
block-adapted absolute basis, spectral synthesis and nuclearity
without that quantitative condition. Identification with a
particular independently defined adelic quotient is a separate
map; MT-M21 begins the comparison with Meyer's construction.
